% Part: first-order-logic
% Chapter: syntax-and-semantics
% Section: terms-formulas

\documentclass[../../../include/open-logic-section]{subfiles}

\begin{document}

\olfileid[pt]{fol}{syn}{frm}

\olsection{Termos e \printtoken{P}{formula}}

Uma vez dada uma linguagem de primeira ordem~$\Lang L$, podemos definir
expressões construídas a partir do vocabulário básico de~$\Lang L$.  Estas
incluem, em particular, \emph{termos} e \emph{!!{formula}s}.

\begin{defn}[Termos]
\ollabel{defn:terms}
O conjunto de \emph{termos}~$\Trm[L]$ de~$\Lang L$ é
definido indutivamente por:
\begin{enumerate}
\item Toda !!{variable} é um termo.
\item Todo !!{constant} de~$\Lang L$ é um termo.
\item Se $f$ é um !!{function} de $n$ lugares e $t_1$, \dots, $t_n$
  são termos, então $\Atom{f}{t_1, \ldots, t_n}$ é um termo.
\tagitem{limitClause}{
  Nada mais é um termo.}{}
\end{enumerate}
Um termo que não contém !!{variable}s é um \emph{termo fechado}.
\end{defn}

\begin{explain}
Os !!{constant}s aparecem em nossa especificação da linguagem e dos
termos como uma categoria separada de símbolos, mas eles poderiam, em vez disso, ter
sido incluídos como !!{function}s de zero lugares.  Poderíamos então dispensar
a segunda cláusula da definição de termos. Bastaria
entender $\Atom{f}{t_1, \ldots, t_n}$ como apenas $f$ sozinho se $n =
0$.
\end{explain}

\begin{defn}[Fórmulas]
\ollabel{defn:formulas}
O conjunto de \emph{!!{formula}s}~$\Frm[L]$ da linguagem~$\Lang L$
é definido indutivamente como segue:
\begin{enumerate}
\tagitem{prvFalse}{$\lfalse$ é uma !!{formula} atômica.}{}

\tagitem{prvTrue}{$\ltrue$ é uma !!{formula} atômica.}{}

\item Se $R$ é um !!{predicate} de $n$ lugares de~$\Lang L$ e $t_1$, \dots,
  $t_n$ são termos de~$\Lang L$, então $\Atom{R}{t_1,\ldots, t_n}$ é uma
  !!{formula} atômica.

\item Se $t_1$ e $t_2$ são termos de~$\Lang L$, então $\Atom{\eq}{t_1, t_2}$
  é uma !!{formula} atômica.

\tagitem{prvNot}{Se $!A$ é !!a{formula}, então $\lnot !A$ é
  !!a{formula}.}{}

\tagitem{prvAnd}{Se $!A$ e $!B$ são !!{formula}s, então $(!A \land
  !B)$ é !!a{formula}.}{}

\tagitem{prvOr}{Se $!A$ e $!B$ são !!{formula}s, então $(!A \lor !B)$
  é !!a{formula}.}{}

\tagitem{prvIf}{Se $!A$ e $!B$ são !!{formula}s, então $(!A \lif !B)$
  é !!a{formula}.}{}

\tagitem{prvIff}{Se $!A$ e $!B$ são !!{formula}s, então $(!A \liff !B)$
  é !!a{formula}.}{}

\tagitem{prvAll}{Se $!A$ é !!a{formula} e $x$ é !!a{variable},
  então $\lforall[x][!A]$ é !!a{formula}.}{}

\tagitem{prvEx}{Se $!A$ é !!a{formula} e $x$ é !!a{variable},
  então $\lexists[x][!A]$ é !!a{formula}.}{}

\tagitem{limitClause}{Nada mais é !!a{formula}.}{}
\end{enumerate}
\end{defn}

\begin{explain}
As definições do conjunto de termos e do conjunto de !!{formula}s são
\emph{definições indutivas}.  Essencialmente, construímos o conjunto de
!!{formula}s em infinitos estágios.  No estágio inicial,
declaramos que todas as fórmulas atômicas são fórmulas; isso corresponde aos
primeiros casos da definição, isto é, os casos para
\iftag{prvTrue}{$\ltrue$, }{}%
\iftag{prvFalse}{$\lfalse$, }{}%
$\Atom{R}{t_1,\dots,t_n}$ e $\Atom{\eq}{t_1,t_2}$. ``!!^{formula}
atômica'' significa, portanto, qualquer !!{formula} desta forma.

Os outros casos da definição dão regras para construir novas
!!{formula}s a partir de !!{formula}s já construídas.  No segundo
estágio, podemos usá-los para construir !!{formula}s a partir de
!!{formula}s atômicas.  No terceiro estágio, construímos novas fórmulas a partir das
fórmulas atômicas e daquelas obtidas no segundo estágio, e assim por diante.  Uma
!!{formula} é qualquer coisa que seja eventualmente construída em tal
estágio, e nada mais.
\end{explain}

Por convenção, escrevemos $\eq$ entre seus argumentos e omitimos os
parênteses: $\eq[t_1][t_2]$ é uma abreviação para
$\Atom{\eq}{t_1,t_2}$.  Além disso, $\lnot \Atom{\eq}{t_1,t_2}$ é
abreviado como $\eq/[t_1][t_2]$. Ao escrever uma fórmula $(!B \ast !C)$
construída a partir de $!B$, $!C$ usando um conectivo de dois lugares~$\ast$,
frequentemente omitiremos o par de parênteses mais externo e escreveremos
simplesmente~$!B \ast !C$.

\begin{intro}
Alguns textos de lógica exigem que a !!{variable}~$x$ deva ocorrer em~$!A$
para que
\iftag{prvEx}{$\lexists[x][!A]$ }{}%
\iftag{notprvEx,notprvAll}{}{e }%
\iftag{prvAll}{$\lforall[x][!A]$ }{}%
contem como
\iftag{notprvEx,notprvAll}{!!a{formula}}{!!{formula}s}.
Nada de ruim acontece se você não exigir isso, e isso torna as coisas
mais fáceis.
\end{intro}

\begin{tagblock}{defNot,defOr,defAnd,defIf,defIff,defTrue,defFalse,defEx,defAll}
\begin{defn}
As fórmulas construídas usando os operadores definidos devem ser entendidas
como segue:

\begin{tagenumerate}{defTrue,defFalse,defNot,defOr,defAnd,defIf,defIff,defEx,defAll}

\tagitem{defTrue}{$\ltrue$ abrevia
  \iftag{prvFalse}{$\lnot\lfalse$}{$(!A \lor \lnot !A)$ para alguma
    !!{formula} atômica fixa~$!A$}.}{}

\tagitem{defFalse}{$\lfalse$ abrevia
  \iftag{prvTrue}{$\lnot\ltrue$}{$(!A \land \lnot !A)$ para alguma
    !!{formula} atômica fixa~$!A$}.}{}

\tagitem{defNot}{$\lnot !A$ abrevia $!A \lif \lfalse$.}{}

\tagitem{defOr}{$!A \lor !B$ abrevia
  \iftag{prvAnd}{$\lnot(\lnot !A \land \lnot !B)$}{$\lnot !A \lif
    !B$}.}{}

\tagitem{defAnd}{$!A \land !B$ abrevia
  \iftag{prvOr}{$\lnot(\lnot !A \lor \lnot !B)$}{$\lnot (!A \lif
    \lnot !B)$}.}{}

\tagitem{defIf}{$!A \lif !B$ abrevia
  \iftag{prvOr}{$\lnot !A \lor !B)$}{$\lnot (!A \land \lnot !B)$}.}{}

\tagitem{defIff}{$!A \liff !B$ abrevia $(!A \lif !B) \land (!B
  \lif !A)$.}{}

\tagitem{defAll}{$\lforall[x][!A]$ abrevia $\lnot\lexists[x][\lnot !A]$.}{}

\tagitem{defEx}{$\lexists[x][!A]$ abrevia $\lnot\lforall[x][\lnot !A]$.}{}
\end{tagenumerate}
\end{defn}
\end{tagblock}

Se trabalhamos em uma linguagem para uma aplicação específica, frequentemente
escreveremos !!{predicate}s e !!{function}s de dois lugares entre os
respectivos termos, por exemplo, $t_1 < t_2$ e $(t_1 + t_2)$ na linguagem
da aritmética e $t_1 \in t_2$ na linguagem da teoria dos conjuntos.  A
função sucessora na linguagem da aritmética é até escrita
convencionalmente \emph{depois} de seu argumento:~$t'$.  Oficialmente, no entanto,
estas são apenas abreviações convencionais para $\Atom{\Obj A^2_0}{t_1, t_2}$,
$\Obj f^2_0(t_1, t_2)$, $\Atom{\Obj A^2_0}{t_1, t_2}$ e $f^1_0(t)$,
respectivamente.

\begin{defn}[Identidade sintática]
O símbolo $\ident$ expressa a identidade sintática entre cadeias de
símbolos, isto é, $!A \ident !B$ sse $!A$ e $!B$ são cadeias de símbolos
do mesmo comprimento e que contêm o mesmo símbolo em cada posição.
\end{defn}

O símbolo $\ident$ pode ser ladeado por cadeias obtidas por
concatenação, por exemplo, $!A \ident (!B \lor !C)$ significa: a cadeia de
símbolos~$!A$ é a mesma cadeia que aquela obtida concatenando
um parêntese de abertura, a cadeia $!B$, o símbolo $\lor$, a
cadeia~$!C$, e um parêntese de fechamento, nesta ordem. Se este é o
caso, então sabemos que o primeiro símbolo de~$!A$ é um parêntese de
abertura, $!A$ contém $!B$ como uma subcadeia (começando no segundo
símbolo), essa subcadeia é seguida por~$\lor$, etc.

Como termos e !!{formula}s são construídos a partir de elementos básicos via definições
indutivas, podemos usar os seguintes princípios de indução para provar
coisas sobre eles.

\begin{lem}[\emph{Princípio de indução sobre termos}]
    \ollabel{lem:trmind}
    Seja $\Lang L$ uma linguagem de primeira ordem.
    Se uma propriedade~$P$ é tal que
    %
    \begin{enumerate}
        \item ela vale para toda !!{variable}~$v$,
        %
        \item ela vale para todo !!{constant}~$a$ de~$\Lang L$, e
        %
        \item ela vale para $f(t_1,\dotsc,t_n)$ sempre que vale para
        $t_1$,~\dots, $t_n$ e $f$~é um !!{function} de $n$ lugares
            de~$\Lang L$
    \end{enumerate}
    (supondo que $t_1$,~\dots, $t_n$ são termos de~$\Lang{L}$),
    então $P$ vale para todo termo em~$\Trm[L]$.
\end{lem}

\begin{prob}
    Prove \olref[fol][syn][frm]{lem:trmind}.
\end{prob}

\begin{lem}[\emph{Princípio de indução sobre !!{formula}s}]
    \ollabel{thm:frmind}
    Seja $\Lang L$ uma linguagem de primeira ordem.
    Se uma propriedade~$P$ vale para todas as !!{formula}s atômicas
    e é tal que
    %
    \begin{enumerate}
        \tagitem{prvNot}{ela vale para $\lnot !A$ sempre que
            vale para~$!A$;}{}
        \tagitem{prvAnd}{ela vale para $(!A \land !B)$
            sempre que vale para $!A$ e~$!B$;}{}
        \tagitem{prvOr}{ela vale para $(!A \lor !B)$
            sempre que vale para $!A$ e~$!B$;}{}
        \tagitem{prvIf}{ela vale para $(!A \lif !B)$
            sempre que vale para $!A$ e~$!B$;}{}
        \tagitem{prvIff}{ela vale para $(!A \liff !B)$
            sempre que vale para $!A$ e~$!B$;}{}
        \tagitem{prvEx}{ela vale para $\lexists[x][!A]$
            sempre que vale para~$!A$;}{}
        \tagitem{prvAll}{ela vale para $\lforall[x][!A]$
            sempre que vale para~$!A$;}{}
  \end{enumerate}
  (supondo que $!A$ e $!B$ são !!{formula}s de~$\Lang{L}$),
  então $P$ vale para todas as fórmulas em~$\Frm[L]$.
\end{lem}

\end{document}
